\section{A uniform overlap bound for the original finite level system}

% BEGIN F74N1
\paragraph{F74N1. Counting the self-dual lattices at each intersection index.}
Let $B$ be the rational alternating form of the given period chart
on $E^r$, of rational dimension $d=[E:\mathbb Q]r$. Its symplectic
basis makes $B$ nondegenerate. Fix one of its self-dual full
$\mathcal O_E$-lattices $L$. Write
$m(L,L')=[L:L\cap L']$. Among all self-dual period lattices there
are at most $m^{2d}$ choices of $L'$ with $m(L,L')=m$ for $m\ge1$.
The only lattice at index one is $L$.

Indeed, forgetting the $\mathcal O_E$-action gives an injection into
the set of full integral lattices self-dual for $B$. The integral
dual condition is the same condition on the underlying additive
group. Alternatively, finite generation and self-duality imply
fullness: the annihilator of the rational span would otherwise contain
a rational line lying in the integral dual and hence in a finitely
generated free abelian group. For a free abelian group of rank $d$,
there are at most $m^d$ subgroups of index $m$. To see this directly,
project onto the last $d-1$ coordinates. If the intersection with
the first coordinate axis is $a\mathbb Z$, the image is a lattice
$H$ of index $m/a$ in $\mathbb Z^{d-1}$. A lift of $H$ is
specified by a homomorphism $H\to\mathbb Z/a$, of which there
are $a^{d-1}$. By induction there are at most
$\sum_{a\mid m}a^{d-1}(m/a)^{d-1}\le m^d$ choices.
The rank-zero case is immediate.

Now choose $M=L\cap L'$, of index $m$ in $L$, and let $M^\vee$
be its integral dual. The inclusion $M\subset L$ gives
$[M^\vee:L]=m$ by transposing an integral basis matrix. Using
$M\subset L'$ gives the same index reversal
$[M^\vee:L']=[L':M]$. Since $[M^\vee:M]=m^2$, both of these
last indices equal $m$. For a given $M$, every possible $L'$ is
therefore an index-$m$ subgroup of $M^\vee$. Applying the preceding
count twice proves the bound; at index one the inclusions force
$L'=L$. The same argument works for any subset of the self-dual
lattices, in particular the original marked level set.
% END F74N1

% BEGIN F74N2
\paragraph{F74N2. The two-lattice inner product in the original finite level.}
Take the source level $D_\ell$ with half ideal $\mathfrak m$,
and finite matrix subgroup $J$. The level assumptions
are $(\mathfrak m,\bar{\mathfrak m})=(\mathfrak m,2)=1$.
Use the actual even column count $n\ge r$, $M>0$ lying in
$\mathfrak m$, a permitted source lattice $\Lambda\in V_{n,M}$,
and the specified normalized input vacua. For the finite test functions
$\Phi_L=D_\ell.\mathrm{phi}(L,\Lambda)$,
\[
 \langle\Phi_L,\Phi_L\rangle=|J|,\qquad
 |\langle\Phi_L,\Phi_{L'}\rangle|
       \le |J|\,[L:L\cap L']^{-n/2}.
\]

Here is the role of the marking in this estimate. Set
$S_L=\mathfrak m\bar{\mathfrak m}^{-1}L$ and
$A_L=\Lambda\otimes S_L$ in the original matrix coordinates.
$S_L$ is the source's ideal quotient since nonzero fractional
ideals are invertible. It is self-dual because its multiplying
ideal has product one with its conjugate. The tensor lattice is
also self-dual for the matrix trace pairing: a local basis of one
factor and its dual reduce this assertion to the duality in the
other. For two levels the common ideal factor cancels from the
local lengths, and $\Lambda$ is locally free of rank $n$.
Taking the sum of lengths at each prime and multiplying the local
indices gives
\[
          [A_L:A_L\cap A_{L'}]=[L:L\cap L']^n.
\]

Put $u_L$ equal to the normalized projection of the original
vacuum onto the $A_L$-vacuum line. For two self-dual tensor
lattices $A,B$, the inner product of their unit vacua has modulus
$[B:A\cap B]^{-1/2}$ with the prescribed common lattice
semicharacter. Indeed average the $B$-translations of the
$A$-vacuum with their inverse characters over $B/(A\cap B)$.
The two characters agree on the intersection. Distinct terms of
this average are orthogonal, as translation by an element of $A$
distinguishes two different cosets by the perfect pairing.
The average is the orthogonal projection onto the $B$-vacuum
line: it is invariant with the prescribed character, and its
inner product with a vector on that line equals the unprojected
inner product. Its squared norm is the reciprocal of the index.
This also proves that the projected vacua being normalized here
are nonzero. The two directions of the relative index coincide by
F74N1's dual-index argument.

In the finite sum defining $\Phi$, the first $r$ rows encode an
element of $J$ in the chosen residual frame and the remaining rows
are zero modulo $\mathfrak m$. Choose one matrix $y_k$ of representatives
for each $k\in J$, simultaneously for $L$ and $L'$. To justify
this, use the lattice with target factor
$2\bar{\mathfrak m}(L\cap L')$. At every prime of $\mathfrak m$
both marked lattices equal the reference lattice, the source
lattice equals the standard one, and $2\bar{\mathfrak m}$ is a
unit locally. Reduction of this full lattice onto the finite
quotient is surjective by the Chinese remainder theorem. The
specified residual frame is a basis there: it spans, and a
surjection between free modules of the same rank over each local
factor is invertible. These residues are therefore distinct.
In the defining sum, a change within a residue class lies in
the denominator lattice with target factor
$2\bar{\mathfrak m}\mathfrak mL$ (or $L'$). Its even lattice
translation fixes the vacuum: write the change as $2z$ with
$z\in A_L$, so its semicharacter is one. A numerator representative
has the form $2y'$; the pairing of $y'$ with $2z$ is integral
by tensor-lattice integrality. Thus the pairing of $2y'$ with $2z$
is even and the Heisenberg product phase is one. Equivalently one may
use the native common representative formula
$\Phi_L=\sum_{k\in J}\lambda(k)\,T_{y_k}u_L$, with exactly the
same $y_k$ for $L'$.

A nonzero cross inner product at $y_k,y_{k'}$ forces
$B(\eta,y_k-y_{k'})\in\mathbb Z$ (the sum of the row pairings)
for every $\eta\in A_L\cap A_{L'}$,
by applying the common translation and cancelling its character.
Thus $y_k-y_{k'}\in(A_L\cap A_{L'})^\vee=A_L+A_{L'}$.
The latter identity follows by dualizing the sum of the two dual
lattices. At the primes of $\mathfrak m$ this sum consists rowwise
of $\mathfrak m$ times the reference lattice. Hence $k=k'$.
For $k=k'$ unitarity of translation and $|\lambda(k)|=1$ leave
$\langle u_L,u_{L'}\rangle$. There are precisely $|J|$ such
terms, proving the two assertions, including the diagonal.
% END F74N2

% BEGIN F74N3
\paragraph{F74N3. The strict bound is uniform in the lattice and in the marking.}
Suppose $n\ge 4d+4$ in F74N2. The original norm field of
$\mathrm{IsSystem}$ now has diagonal $|J|>0$, majorant
$b(L,L')=[L:L\cap L']^{-n/2}$ and off-diagonal sum at most $0.65$.
Indeed for each finite set of lattices other than $L$, group by the
integer $m\ge2$ of F74N1. The sum of the majorant is bounded by
\[
   \sum_{m=2}^{\infty} m^{2d-n/2}
       \le \sum_{m=2}^{\infty}m^{-2}
       =\pi^2/6-1<0.65 .
\]
All terms are nonnegative, so this proves summability and the same
bound for the entire set. The argument takes the supremum over
finite sums and needs no ordering of the period lattices. The
constant is independent of the class representative $\Lambda$
and of $L$. In particular the diagonal term dominates the entire
off-diagonal sum even at the endpoint $n=4d+4$.

For the finite positive class weights of ICL12 put
$a=|J|\sum_{[\Lambda]}\mathrm{wt}(M,\Lambda)$.
The weighted matrix has diagonal $a$ and the off-diagonal row sum
is at most $0.65a$. On the source space with its prescribed
covariance, once the level action preserves that space, its inverse
is consequently given by the Neumann series with norm at most
$1/(0.35a)$. This supplies the quantitative norm hypothesis
previously assumed in ICL12 for these test vectors. The covariance,
support, cuspidality and doubled-family admissibility fields of
the full level system have independent content.
% END F74N3
